\documentclass[10pt]{article}
\usepackage[T1]{fontenc}
\usepackage[utf8]{inputenc}
\usepackage[french]{babel}
\usepackage{amsmath,amssymb,mathrsfs,mathpazo}
\usepackage{geometry}
\geometry{a4paper,margin=2cm}
\usepackage{graphicx,array,tabularx,tabularray,colortbl,multirow,multicol,booktabs,enumitem,pifont,ulem,fancybox,diagbox,cancel}
\usepackage{pstricks,pstricks-add,pst-plot,pst-tree,pst-node,pst-text,pst-func,pst-math,pst-bspline,pst-3dplot,pst-eucl,pst-solides3d}
\let\pscalculate\undefined
\usepackage{pst-bezier}
\usepackage{tikz,pgfplots,tkz-tab}
\usetikzlibrary{arrows,arrows.meta,patterns,calc,positioning,trees,decorations.markings,shapes.geometric,shapes.misc,shapes.arrows,intersections,decorations.pathmorphing,decorations.pathreplacing,backgrounds,shadows,babel}
\pgfplotsset{compat=1.18}
\usepgfplotslibrary{fillbetween}
\usepackage[np]{numprint}
\usepackage{listings,makecell,pas-tableur,twoopt}
\let\barre\undefined
\usepackage[pstricks]{tablvar}
\definecolor{darkgreen}{RGB}{0,100,0}
\providecommand{\red}{\color{red}}
\providecommand{\blue}{\color{blue}}
\providecommand{\green}{\color{green}}
\providecommand{\black}{\color{black}}
\providecommand{\white}{\color{white}}
\providecommand{\gray}{\color{gray}}
\providecommand{\cyan}{\color{cyan}}
\providecommand{\magenta}{\color{magenta}}
\providecommand{\yellow}{\color{yellow}}
\providecommand{\violet}{\color{violet}}
\providecommand{\orange}{\color{orange}}
\providecommand{\euro}{\texteuro}
\providecommand{\EUR}{\texteuro}
\providecommand{\ieme}{\textsuperscript{e}}
\providecommand{\degres}{\ensuremath{^\circ}}
\providecommand{\degree}{\ensuremath{^\circ}}
\providecommand{\vv}[1]{\overrightarrow{#1}}
\providecommand{\Coord}[1]{#1}
\providecommand{\decosix}{}
\providecommand{\starredbullet}{\ensuremath{\bullet}}
\pagestyle{empty}
\setlength{\parindent}{0pt}
\usepackage{adjustbox}\DeclareUnicodeCharacter{2236}{:}\DeclareUnicodeCharacter{2190}{\ensuremath{\leftarrow}}
\providecommand{\R}{\mathbb{R}}
\providecommand{\N}{\mathbb{N}}
\providecommand{\Z}{\mathbb{Z}}
\providecommand{\Q}{\mathbb{Q}}
\providecommand{\D}{\mathbb{D}}
\providecommand{\C}{\mathbb{C}}
\providecommand{\e}{{\mathrm{e}}}
\providecommand{\vect}[1]{\overrightarrow{#1}}
\providecommand{\vectt}[1]{\overrightarrow{\mathrm{#1}}}
\providecommand{\vvt}[1]{\overrightarrow{\mathrm{#1}}}
\providecommand{\np}[2][]{{#2~#1}}
\providecommand{\vv}[1]{{\overrightarrow{#1}}}
\renewcommand{\d}{{\mathrm{d}}}
\providecommand{\degres}{{^\circ}}
\providecommand{\degre}{{^\circ}}
\providecommand{\blue}{\color{blue}}
\providecommand{\red}{\color{red}}
\providecommand{\enskip}{\quad}
\providecommand{\llbracket}{\lbrack\!\lbrack}
\providecommand{\rrbracket}{\rbrack\!\rbrack}
\providecommand{\euro}{{\text{€}}}
\providecommand{\no}{{\text{n°}}}
\providecommand{\interval}[2]{{[#1;#2]}}
\providecommand{\intervFO}[2]{{[#1;#2[}}
\providecommand{\dotfill}{\ldots}

\providecommand{\ds}{\displaystyle}
\providecommand{\textding}[1]{#1}
\providecommand{\Oij}{\(\mathrm{O};\vec{\imath},\vec{\jmath}\)}
\providecommand{\Oijk}{\(\mathrm{O};\vec{\imath},\vec{\jmath},\vec{k}\)}
\providecommand{\Ouv}{\(\mathrm{O};\vec{u},\vec{v}\)}
\providecommand{\intervFF}[2]{\ensuremath{[#1;#2]}}
\renewcommand{\euro}{\ifmmode\text{\texteuro}\else\texteuro\fi}
\DeclareUnicodeCharacter{00B0}{\ifmmode^{\circ}\else\textdegree\fi}
\DeclareUnicodeCharacter{20AC}{\ifmmode\text{\texteuro}\else\texteuro\fi}
% Énoncés : Éducation nationale. Transcription et corrigés : APMEP, Denis Vergès et contributeurs. Découpage et classement : CoopMaths, Sébastien Lozano, Éric Elter et contributeurs. Adaptation de la présentation : Maths Explorer.
% CC BY-SA : https://coopmaths.fr/www/about/
% https://coopmaths.fr/alea/static/bac/2026/tex/bac_2026_06_sujet2_polynesie_3_cor.tex
\begin{document}
\noindent\textit{APMEP / CoopMaths. CC BY-SA. Présentation adaptée par Maths Explorer.}\par\medskip
\begin{enumerate}
	\item \begin{enumerate}
		\item On a $\vect{\text{AB}}\begin{pmatrix}2\\-7\\-2\end{pmatrix}$ et $\vect{\text{AC}}\begin{pmatrix}-4\\1\\2\end{pmatrix}$. On voit que les coordonnées de ces deux vecteurs ne sont pas proportionnelles : les vecteurs ne sont pas colinéaires donc les points A, B et C ne sont pas alignés : ils déterminent donc un plan (ABC).

		\item On a $\vect{n} \cdot \vect{\text{AB}} = 12 + 14 - 26 = 0$ ; de même

		$\vect{n} \cdot \vect{\text{AB}} = - 24 - 2 + 26 = 0$.

		Le vecteur $\vect{n}$ est orthogonal à deux vecteurs non colinéaires du plan (ABC) : c'est donc un vecteur normal à ce plan (ABC).

		\item On sait que les coordonnées du vecteur $\vect{n}$ sont les coefficients respectifs de $x, \: y$ et $z$ dans l'équation du plan (ABC), soit

		$M(x ; y ; z) \in (\text{ABC}) \iff 6x - 2y + 13z + d = 0$,

		avec $d \in \R$.

		En particulier on a $A(-1 ; 1 ; 1) \in \text{ABC}) \iff 6\times (-1) - 2\times 1 + 13\times 1 + d = 0 \iff - 6 - 2 + 13 + d = 0 \iff 5 + d = 0 \iff d = - 5$.

		Conclusion : \quad $\mathrm{M}(x ; y ; z) \in (\text{ABC}) \iff 6x - 2y + 13z - 5 = 0$.
	\end{enumerate}

	\item Comme on l'a vu à la question précédente les coefficients respectifs de $x$, $y$ et $z$ dans l'équation du plan $(P)$ sont les coordonnées d'un vecteur $\vect{p}$ normal à ce plan.

	On a donc $\vect{p}\begin{pmatrix}8\\- 2\\-4\end{pmatrix}$.

	Or $\vect{n} \cdot \vect{p} = 6 \times 8 - 2 \times (- 2) + 13 \times (- 4) = 48 + 4 - 52 = 0$ : d'après la propriété admise les deux vecteurs normaux aux deux plans sont orthogonaux donc les deux plans sont orthogonaux.

	\item La droite $(d)$ est sécante au plan (ABC) s'il existe un point dont les coordonnées vérifient les équations paramétriques de $(d)$ et l'équation cartésienne du plan (ABC), autrement dit si les coordonnées $(x ; y ; z)$ vérifient le système d'équations :
		\[\left\{\begin{array}{l c l}
			x&=&- 4 + t\\
			y&=&- 11 + 6t\\
			z&=&2 - t\\
			6x - 2y + 13z - 5 &=& 0
		\end{array}\right.\]

		En remplaçant $x$, $y$ et $z$ par leurs expressions en fonction de $t$ dans la dernière équation, on obtient :

		$6(- 4 + t)- 2(- 11 + 6t) + 13(2 - t) - 5 = 0 \iff - 24 + 6t + 22 - 12t + 26 - 13t - 5 = 0$,

		soit enfin \quad $19 - 19t = 0 \iff 19 = 19t \iff 1 = t$.

		En reportant cette valeur $t = 1$ dans les équations paramétriques de $(d)$, on obtient \quad $x = - 4 + 1 = - 3$, \: $y = - 11 + 6 = - 5$, \:$z = 2 - 1 = 1$.

		La droite $(d)$ et le plan (ABC) ont bien un point commun E$(- 3 ; -5 ; 1)$.

		\item \begin{enumerate}
			\item H est le projeté orthogonal de A sur $(d)$ si :

			\begin{itemize}
				\item  H $\in (d)$, ce qui est aisé à voir en prenant la valeur $t = 2$ dans les équations paramétriques de $(d)$ qui conduit bien aux coordonnées $(- 2 ; 1 ; 0)$ qui sont bien celles de H;
				\item la droite (AH) est perpendiculaire à $(d)$ ;
			\end{itemize}

			Or $\vect{\text{AH}}\begin{pmatrix}-1\\0\\- 1\end{pmatrix}$ et $\vect{\delta}\begin{pmatrix}1\\6\\- 1\end{pmatrix}$ est un vecteur directeur de la droite $(d)$ (d'après ses équations paramétriques)

			et $\vect{\text{AH}} \cdot \vect{\delta} = - 1 + 0 + 1 = 0$, ce qui montre que les vecteurs $\vect{\text{AH}}$ et $\vect{\delta}$ sont orthogonaux donc que les droites (AH) et ($d$) sont perpendiculaires.

			Le point H est le projeté orthogonal de A sur $(d)$.

			\item La distance du point A à la droite $(d)$ est la distance de A à son projeté orthogonal sur la droite $(d)$, soit la distance AH.

			Des coordonnées du vecteur $\vect{\text{AH}}$, on calcule AH$^2 = (- 1)^2 + 0^2 + 1^2 = 1 + 1 = 2$. Conclusion AH $ = \sqrt 2$.

			\item H et E sont deux points de $(d)$ et la droite (AH) est perpendiculaire à $(d)$ : le triangle AHE est donc rectangle en H et son aire est égale au demi-produit des longueurs de ses côtés perpendiculaires :

			$\mathcal{A}(\text{AHE}) = \dfrac{\text{AH} \times \text{HE}}{2}$

			On a HE$^2 = (- 3 + 2)^2 + (- 5 - 1)^2 + (1 - 0)^2 = 1 + 36 + 1 = 38$, d'où HE $ = \sqrt{38}$ et

			$\mathcal{A}(\text{AHE}) = \dfrac{\sqrt 2 \times \sqrt{38}}{2} = \dfrac{\sqrt{76}}{2} = \dfrac{\sqrt{76}}{\sqrt4} = \sqrt{\dfrac{76}{4}} = \sqrt{19}$.
		\end{enumerate}
	\end{enumerate}

\end{document}
